% Options for packages loaded elsewhere
\PassOptionsToPackage{unicode}{hyperref}
\PassOptionsToPackage{hyphens}{url}
%
\documentclass[
]{article}
\usepackage{amsmath,amssymb}
\usepackage{iftex}
\ifPDFTeX
  \usepackage[T1]{fontenc}
  \usepackage[utf8]{inputenc}
  \usepackage{textcomp} % provide euro and other symbols
\else % if luatex or xetex
  \usepackage{unicode-math} % this also loads fontspec
  \defaultfontfeatures{Scale=MatchLowercase}
  \defaultfontfeatures[\rmfamily]{Ligatures=TeX,Scale=1}
\fi
\usepackage{lmodern}
\ifPDFTeX\else
  % xetex/luatex font selection
\fi
% Use upquote if available, for straight quotes in verbatim environments
\IfFileExists{upquote.sty}{\usepackage{upquote}}{}
\IfFileExists{microtype.sty}{% use microtype if available
  \usepackage[]{microtype}
  \UseMicrotypeSet[protrusion]{basicmath} % disable protrusion for tt fonts
}{}
\makeatletter
\@ifundefined{KOMAClassName}{% if non-KOMA class
  \IfFileExists{parskip.sty}{%
    \usepackage{parskip}
  }{% else
    \setlength{\parindent}{0pt}
    \setlength{\parskip}{6pt plus 2pt minus 1pt}}
}{% if KOMA class
  \KOMAoptions{parskip=half}}
\makeatother
\usepackage{xcolor}
\usepackage[margin=1cm]{geometry}
\usepackage{graphicx}
\makeatletter
\def\maxwidth{\ifdim\Gin@nat@width>\linewidth\linewidth\else\Gin@nat@width\fi}
\def\maxheight{\ifdim\Gin@nat@height>\textheight\textheight\else\Gin@nat@height\fi}
\makeatother
% Scale images if necessary, so that they will not overflow the page
% margins by default, and it is still possible to overwrite the defaults
% using explicit options in \includegraphics[width, height, ...]{}
\setkeys{Gin}{width=\maxwidth,height=\maxheight,keepaspectratio}
% Set default figure placement to htbp
\makeatletter
\def\fps@figure{htbp}
\makeatother
\setlength{\emergencystretch}{3em} % prevent overfull lines
\providecommand{\tightlist}{%
  \setlength{\itemsep}{0pt}\setlength{\parskip}{0pt}}
\setcounter{secnumdepth}{-\maxdimen} % remove section numbering
\usepackage{amsmath}
\usepackage{mathtools}
\usepackage{amsthm}
\usepackage{color}
\usepackage{xcolor}
\usepackage{geometry}
\usepackage{enumerate}
\ifLuaTeX
  \usepackage{selnolig}  % disable illegal ligatures
\fi
\IfFileExists{bookmark.sty}{\usepackage{bookmark}}{\usepackage{hyperref}}
\IfFileExists{xurl.sty}{\usepackage{xurl}}{} % add URL line breaks if available
\urlstyle{same}
\hypersetup{
  pdftitle={Regression Analysis - STAT 482 - Exam 1: Due December 2, 2021},
  pdfauthor={Alex Towell (atowell@siue.edu)},
  hidelinks,
  pdfcreator={LaTeX via pandoc}}

\title{Regression Analysis - STAT 482 - Exam 1: Due December 2, 2021}
\author{Alex Towell
(\href{mailto:atowell@siue.edu}{\nolinkurl{atowell@siue.edu}})}
\date{}

\begin{document}
\maketitle

\newcommand{\sos}[1]{\mathrm{SS_{#1}}}
\newcommand{\ms}[1]{\mathrm{MS_{#1}}}
\newcommand{\sd}{\operatorname{sd}}
\newcommand{\var}{\operatorname{var}}
\newcommand{\expect}{\operatorname{E}}
\newcommand{\corr}{\operatorname{cor}}
\newcommand{\cov}{\operatorname{cov}}
\newcommand{\se}{\operatorname{se}}
\newcommand{\eval}[2]{\left. #1 \right\vert_{#2}}
\newcommand{\degf}[1]{\mathrm{df_{#1}}}
\newcommand{\entropy}{\operatorname{H}}
\newcommand{\param}[1]{\textcolor{blue}{\texttt{#1}}}
\newcommand{\ci}{\operatorname{CI}}

\hypertarget{problem-1}{%
\section{Problem 1}\label{problem-1}}

\begin{quote}
Experience with a certain type of plastic indicates that a relationship
exists between the hardness (measured in Brinell units) of items molded
from the plastic (response variable \(y\) ) and the elapsed time
(measured in hours) since the termination of the molding process (input
variable \(x\)). Sixteen batches of the plastic were made, and from each
batch one test item was molded. Each test item was randomly assigned to
one of the four predetermined time levels. The data is available on
Blackboard as a csv file.
\end{quote}

\vspace{50pt}

\hypertarget{a}{%
\subsection{(a)}\label{a}}

\begin{quote}
State the simple linear regression model.
\end{quote}

\vspace{50pt}

\hypertarget{b}{%
\subsection{(b)}\label{b}}

\begin{quote}
Provide an interpretation for the regression parameter
\end{quote}

\vspace{50pt}

\hypertarget{c}{%
\subsection{(c)}\label{c}}

\begin{quote}
Compute an interval estimate for \(\beta_1\), stated in the context of
the problem.
\end{quote}

\vspace{50pt}

\hypertarget{d}{%
\subsection{(d)}\label{d}}

\begin{quote}
Compute a confidence interval for \(h\) at \(x_h = 40\) hours.
\end{quote}

\vspace{50pt}

\hypertarget{e}{%
\subsection{(e)}\label{e}}

\begin{quote}
Compute a prediction interval for \(Y_{h(\text{new})}\) at \(x_h = 40\)
hours.
\end{quote}

\vspace{50pt}

\hypertarget{f}{%
\subsection{(f)}\label{f}}

\begin{quote}
Explain the di¤erence between a confidence interval and a prediction
interval, stated in the context of the problem.
\end{quote}

\vspace{50pt}

\hypertarget{g}{%
\subsection{(g)}\label{g}}

\begin{quote}
State the equations for \(E(\ms{R})\) and \(\ms{E}\).
\end{quote}

\vspace{50pt}

\hypertarget{h}{%
\subsection{(h)}\label{h}}

\begin{quote}
Use part (g) to explain a motivation behind the F test for input
effects.
\end{quote}

\vspace{50pt}

\hypertarget{i}{%
\subsection{(i)}\label{i}}

\begin{quote}
Compute the \(F^*\) statistic and the \(p\)-value. Provide an
interpretation of the result, stated in the context of the problem.
\end{quote}

\vspace{50pt}

\hypertarget{j}{%
\subsection{(j)}\label{j}}

\begin{quote}
Compute the coefficient of determination \(r^2\). Provide an
interpretation of the result, stated in the context of the problem.
\end{quote}

\vspace{50pt}

\hypertarget{k}{%
\subsection{(k)}\label{k}}

\begin{quote}
State the full model (saturated model) in testing for regression model
fit.
\end{quote}

\vspace{50pt}

\hypertarget{l}{%
\subsection{(l)}\label{l}}

\begin{quote}
Compute \(\sos{PE}\) and \(\sos{LF}\) for testing the fit of the linear
regression model.
\end{quote}

\vspace{50pt}

\hypertarget{m}{%
\subsection{(m)}\label{m}}

\begin{quote}
Compute the \(F_{LF}\) statistic and the \(p\)-value. Provide an
interpretation of the result, stated in the context of the problem.
\end{quote}

\vspace{50pt}

\hypertarget{n}{%
\subsection{(n)}\label{n}}

\begin{quote}
Create a plot comparing the fitted values from the regression model with
the fitted values from the saturated model.
\end{quote}

\vspace{50pt}

\hypertarget{o}{%
\subsection{(o)}\label{o}}

\begin{quote}
Explain an advantage of using the regression estimate \(\hat{\mu}_{40}\)
instead of the sample mean \(\bar{y}_{40}\) and explain when this
advantage will lead to a more accurate estimator.
\end{quote}

\vspace{50pt}

\hypertarget{problem-2}{%
\section{Problem 2}\label{problem-2}}

\begin{quote}
Measurements were made on men involved in a physical fitness course. The
input variables under consideration are age (in years), weight (in kgs),
time to run 1.5 miles (in minutes), heart rate while resting (in bpm),
heart rate while running, and maximum heart rate. The goal is to
determine which input variables are needed to model the oxygen uptake
rate (ml/kg body weight per minute). The data is available on Blackboard
as a csv file.
\end{quote}

\vspace{50pt}

\hypertarget{a-1}{%
\subsection{(a)}\label{a-1}}

\begin{quote}
Describe the goal of the discrepancy function approach to model
selection. How is the best model defined? What are the two sources of
model error?
\end{quote}

\vspace{50pt}

\hypertarget{b-1}{%
\subsection{(b)}\label{b-1}}

\begin{quote}
Plot the \(C_p\) statistic against the model dimension. Plot the \(C_p\)
statistic against the candidate models. Which variables are included in
the selected model?
\end{quote}

\vspace{50pt}

\hypertarget{c-1}{%
\subsection{(c)}\label{c-1}}

\begin{quote}
Test the selected model against its best competitor having fewer
parameters. (Compute the test statistic and the \(p\)-value.) How does
discrepancy based model selection compare to \(p\)-value based
selection?
\end{quote}

\end{document}
